% --- Шрифты (XeLaTeX) ---
\usepackage{iftex}
\usepackage{fontspec}

\ifxetex
  \IfFontExistsTF{Times New Roman}{
    \setmainfont{Times New Roman}
    \setsansfont{Times New Roman}
    \setmonofont{Courier New}
  }{
    % Liberation — метрически совместим с Times, есть кириллица,
    % ставится в CI через extra_system_packages: "font-liberation"
    \setmainfont{Liberation Serif}
    \setsansfont{Liberation Sans}
    \setmonofont{Liberation Mono}
  }
\else
  \PackageError{preamble}{This document requires XeLaTeX}{}
\fi
% --- Графика и TikZ ---
\usepackage{graphicx}
\usepackage{tikz}
\usetikzlibrary{arrows.meta, positioning, shapes, calc}
\usepackage[most]{tcolorbox}   % Для красивых блоков
\usepackage{fontawesome5}       % Для иконок (работает с XeLaTeX)

% --- Стиль для "Внимание" ---
\newtcolorbox{warningbox}[1][]{
    enhanced,
    colback=red!5!white,      % Светло-красный фон
    colframe=red!75!black,    % Темно-красная рамка
    fonttitle=\bfseries,
    title={\faExclamationTriangle~Внимание}, % Иконка + заголовок
    #1
}

% --- Стиль для "На заметку" ---
\newtcolorbox{notebox}[1][]{
    enhanced,
    colback=blue!5!white,     % Светло-синий фон
    colframe=blue!60!black,   % Синяя рамка
    fonttitle=\bfseries,
    title={\faPen~На заметку}, % Иконка ручки
    #1
}

%--- Стиль для "Внимание" 2й вариант
\newtcolorbox{warningbox2}[1][]{
    enhanced,
    colback=red!5!white,
    colframe=red!75!black,
    overlay={
        \node[anchor=north west, inner sep=10pt] at (frame.north west) 
        {\Huge\color{red!75!black}\faExclamationTriangle};
    },
    left=1.5cm, % Сдвигаем текст вправо, чтобы не наезжал на иконку
    #1
}

%--- Стиль для "На заметку" 2й вариант
\newtcolorbox{notebox2}[1][]{
    enhanced,
    colback=blue!5!white,
    colframe=blue!60!black,,
    overlay={
        \node[anchor=north west, inner sep=10pt] at (frame.north west) 
        {\Huge\color{blue!60!black}\faPen};
    },
    left=1.5cm, % Сдвигаем текст вправо, чтобы не наезжал на иконку
    #1
}

% --- Язык ---
\usepackage[english,russian]{babel}

% Явная фиксация (обычно не требуется)
\makeatletter
\renewcommand{\russianhyphenmins}{22}
\makeatother

% --- Математика ---
\usepackage{amsmath,amssymb}

% --- Код ---
\usepackage{listings}
\usepackage{xcolor}

% --- Геометрия страницы ---
\usepackage{geometry}
\geometry{margin=2.5cm}

% --- Подписи к рисункам ---
\usepackage{caption}
\captionsetup[figure]{
  name=Рис.,
  labelsep=endash,
  justification=centering,
  font=small,
  labelfont=bf,
  skip=6pt
}

%
\addto\captionsrussian{%
  \renewcommand{\figurename}{Рис.}%
  \renewcommand{\tablename}{Таблица}%
  \renewcommand{\lstlistingname}{Листинг}%
}


% --- Плавающие объекты ---
\usepackage{float}

% --- Типографика: убиваем badbox'ы ---
\usepackage{microtype}
\usepackage{xurl}
\emergencystretch=1em
\Urlmuskip=0mu plus 1mu\relax

% --- Настройка листингов ---
\lstset{
    basicstyle=\ttfamily\small,
    breaklines=true,
    frame=single,
    language=Python,
    showstringspaces=false,
    keywordstyle=\color{blue},
    commentstyle=\color{gray},
    stringstyle=\color{red},
    extendedchars=true,
    inputencoding=utf8,
    literate=%
      {а}{{\cyra}}1 {б}{{\cyrb}}1 {в}{{\cyrv}}1 {г}{{\cyrg}}1
      {д}{{\cyrd}}1 {е}{{\cyre}}1 {ё}{{\cyryo}}1 {ж}{{\cyrzh}}1
      {з}{{\cyrz}}1 {и}{{\cyri}}1 {й}{{\cyrishrt}}1 {к}{{\cyrk}}1
      {л}{{\cyrl}}1 {м}{{\cyrm}}1 {н}{{\cyrn}}1 {о}{{\cyro}}1
      {п}{{\cyrp}}1 {р}{{\cyrr}}1 {с}{{\cyrs}}1 {т}{{\cyrt}}1
      {у}{{\cyru}}1 {ф}{{\cyrf}}1 {х}{{\cyrh}}1 {ц}{{\cyrc}}1
      {ч}{{\cyrch}}1 {ш}{{\cyrsh}}1 {щ}{{\cyrshch}}1 {ъ}{{\cyrhrdsn}}1
      {ы}{{\cyrery}}1 {ь}{{\cyrsftsn}}1 {э}{{\cyrerev}}1 {ю}{{\cyryu}}1
      {я}{{\cyrya}}1
      {А}{{\CYRA}}1 {Б}{{\CYRB}}1 {В}{{\CYRV}}1 {Г}{{\CYRG}}1
      {Д}{{\CYRD}}1 {Е}{{\CYRE}}1 {Ё}{{\CYRYO}}1 {Ж}{{\CYRZH}}1
      {З}{{\CYRZ}}1 {И}{{\CYRI}}1 {Й}{{\CYRISHRT}}1 {К}{{\CYRK}}1
      {Л}{{\CYRL}}1 {М}{{\CYRM}}1 {Н}{{\CYRN}}1 {О}{{\CYRO}}1
      {П}{{\CYRP}}1 {Р}{{\CYRR}}1 {С}{{\CYRS}}1 {Т}{{\CYRT}}1
      {У}{{\CYRU}}1 {Ф}{{\CYRF}}1 {Х}{{\CYRH}}1 {Ц}{{\CYRC}}1
      {Ч}{{\CYRCH}}1 {Ш}{{\CYRSH}}1 {Щ}{{\CYRSHCH}}1 {Ъ}{{\CYRHRDSN}}1
      {Ы}{{\CYRERY}}1 {Ь}{{\CYRSFTSN}}1 {Э}{{\CYREREV}}1 {Ю}{{\CYRYU}}1
      {Я}{{\CYRYA}}1
}

% --- Гиперссылки (последним) ---
\usepackage[hidelinks]{hyperref}